% Generated by `taskgen course examples` from number_theory.primitive_root (seed 7).
% Do not edit: change examples.yaml and regenerate.
\begin{example}[модуль без первообразных корней]\label{ex:primroot-none}
Найти наименьший первообразный корень по модулю $n = 76$ или доказать, что первообразных корней по модулю $76$ нет.

\par\noindent\textit{Решение.} $n = 76 = 2^{2} \cdot 19$, $\varphi(n) = 36$. Первообразные корни существуют только по модулям $1, 2, 4, p^k, 2p^k$ ($p$ нечётное простое), а $76$ не такой. Докажем, что порядок любого $x \in (\mathbb{Z}/76\mathbb{Z})^*$ меньше $\varphi(n)$.\par Разложим $n = a b = 4 \cdot 19$, $\gcd(a, b) = 1$, $a, b > 2$; тогда $\varphi(a) = 2$ и $\varphi(b) = 18$ чётны, и для $\gcd(x, n) = 1$ \[\begin{gathered}x^{\varphi(n)/2} = \bigl(x^{\varphi(a)}\bigr)^{\varphi(b)/2} \equiv 1 \pmod{a},\\ x^{\varphi(n)/2} = \bigl(x^{\varphi(b)}\bigr)^{\varphi(a)/2} \equiv 1 \pmod{b},\end{gathered}\] т.\,е. $x^{18} \equiv 1 \pmod{76}$ и $\ord(x) \le 18 < 36$.\par Точнее, максимальный порядок равен функции Кармайкла $\lambda(n) = \operatorname{lcm}(\lambda(2^{2}), \lambda(19)) = \operatorname{lcm}(2, 18) = 18 < 36$.\[\text{первообразных корней по модулю } 76 \text{ нет}.\]
\end{example}
